% BEGIN REV09_NIVELES_2_5
\subsection{Compatibilité des lecteurs modulaires de niveaux deux et cinq}
\label{corpus9:iv:niveles}

Sur la réalisation modulaire postérieure à la construction HMT,
fixons \(\operatorname{Im}\tau>0\), \(q=e^{2\pi i\tau}\) et
\(q^{1/5}=e^{2\pi i\tau/5}\). Le lecteur de Rogers--Ramanujan est

\[
 R(q)=\frac{q^{1/5}}{1+\dfrac{q}{1+\dfrac{q^2}{1+\cdots}}},
 \qquad x=R(q)^5.
\]
L’identité modulaire de cette réalisation, avec la normalisation
de \(j\) utilisée dans \eqref{km:j-nivel2}, est
\begin{equation}
 j=-\frac{(x^4-228x^3+494x^2+228x+1)^3}
          {x(x^2+11x-1)^5}.
 \label{corpus9:iv:nivel5-identidad}
\end{equation}
L’identité constitue la prémisse de réalisation modulaire du
transport algébrique suivant ; les germes, l’orientation et
les représentations d’auto-échelle restent ceux de APP--TRIT--TPK.

\begin{proposition}[Quotient pentadique et sélection de branche]
\label{corpus9:iv:nivel5-cociente}
Pour \(x\ne0\), soit \(u=x-x^{-1}\). Dans le domaine \(u\ne-11\),
\begin{equation}
 j=-\frac{(u^2-228u+496)^3}{(u+11)^5}.
 \label{corpus9:iv:nivel5-u}
\end{equation}
L’involution \(x\mapsto-1/x\) conserve \(u\). Ses préimages sont
\[
 x=\frac{u\pm\sqrt{u^2+4}}2.
\]
Le revêtement quadratique se ramifie en \(u=\pm2i\) ; l’application
rationnelle ultérieure \(u\mapsto j\) est de degré six.
\end{proposition}
\begin{proof}
Les factorisations exactes sont
\[
 \begin{aligned}
 x^4-228x^3+494x^2+228x+1
   &=x^2(u^2-228u+496),\\
 x^2+11x-1&=x(u+11).
 \end{aligned}
\]
Leur substitution dans \eqref{corpus9:iv:nivel5-identidad} prouve
\eqref{corpus9:iv:nivel5-u}. L’équation de la fibre est
\(x^2-ux-1=0\) ; son discriminant donne les points de ramification.
Le numérateur de \eqref{corpus9:iv:nivel5-u} est de degré six
et le dénominateur de degré cinq. En \(u=-11\), le polynôme
du numérateur vaut \(3125\), de sorte qu’ils n’ont pas de
facteur commun. Cela prouve le degré indiqué.
\end{proof}

Pour \(u\) réel, les deux racines ont des signes opposés. Dans le
domaine complexe, la branche choisie dans le paramétrage
par \(\tau\) est conservée. La lecture \(j\) seule conserve une fibre
génériquement de six valeurs de \(u\).


\begin{proposition}[Limite d’auto-échelle et compatibilité des niveaux]
\label{corpus9:iv:niveles-ramas}
Dans la limite réelle \(q\to1^-\),
\[
 R(q)\longrightarrow\varphi_{\mathrm{HMT}}^{-1},\qquad
 x\longrightarrow\varphi_{\mathrm{HMT}}^{-5}
     =\frac{5\sqrt5-11}{2},\qquad u\longrightarrow-11.
\]
Lorsque \(t=t_2(\tau)\) et \(u\) sont lus sur le même paramètre
modulaire, ils satisfont
\begin{equation}
 (t+256)^3(u+11)^5
   +t^2(u^2-228u+496)^3=0.
 \label{corpus9:iv:niveles-enlace}
\end{equation}
En posant \(\epsilon=u+11\), les branches qui satisfont les limites
indiquées ont les lois
\begin{align}
 \epsilon\to0,\ t\to0:
 &\quad \frac{t^2}{(-\epsilon)^5}\longrightarrow
          \frac{2^{24}}{5^{15}},\\
 \epsilon\to0,\ t\to\infty:
 &\quad t\epsilon^5\longrightarrow-5^{15}.
 \label{corpus9:iv:niveles-asintotica}
\end{align}
\end{proposition}
\begin{proof}
Pour \(0<q<1\), les réduites positives alternées de la
fraction continue encadrent \(R(q)\). Chaque réduite fixée
tend, lorsque \(q\to1^-\), vers la réduite correspondante de
la fraction continue dont tous les numérateurs sont égaux à un.
Les deux sous-suites de ces réduites ont pour limite
\(\varphi_{\mathrm{HMT}}^{-1}\), dans sa réalisation quadratique
déjà établie. L’encadrement prouve la première limite. Élever à
la puissance cinq et utiliser \(\varphi^2=\varphi+1\) donne les
deux autres.


Égaler \eqref{km:j-nivel2} et \eqref{corpus9:iv:nivel5-u}
prouve \eqref{corpus9:iv:niveles-enlace}. La substitution
\(u=\epsilon-11\) donne
\[
 \frac{(t+256)^3}{t^2}
   =-\frac{(3125-250\epsilon+\epsilon^2)^3}{\epsilon^5}.
\]
Dans la branche \(t\to0\), l’identité se réarrange sous la forme
\[
 \frac{t^2}{(-\epsilon)^5}
   =\frac{(t+256)^3}{(3125-250\epsilon+\epsilon^2)^3}.
\]
La limite est \(256^3/3125^3=2^{24}/5^{15}\).
Sur la branche \(t\to\infty\), multiplier par \(\epsilon^5\)
et écrire le membre de gauche comme
\(t\epsilon^5(1+256/t)^3\) donne la seconde limite.
\end{proof}

La relation conserve les branches d’une correspondance sur le
même \(j\). Le choix de \(\tau\), ou du relèvement
correspondant, fixe la paire transportée. Le pôle
pentadique et l’inversion de Fricke maintiennent leurs domaines et leurs
opérations respectifs.
% END REV09_NIVELES_2_5
